\section{Rekonstruktion bei zwei Produktwerten}
\hypertarget{nullproducts.main}{}

Nullprodukte erfassen, welche Faktoren einander zu Null multiplizieren.
Für eine Halbgruppe mit genau zwei Produktwerten bestimmen sie bereits,
welcher der beiden Werte bei jedem Produkt auftritt. Die Schwierigkeit ist,
diese Information aus der abstrakten Potenzhalbgruppe zurückzugewinnen:
Ein Isomorphismus der Potenzhalbgruppen muss Einermengen nicht erhalten.

Die allgemeinen Grundlagen dazu stehen in
\FormulaRefAuto{NullProfileTwoPointReconstruction}[thm] in Band~33.
Sie rekonstruieren die Häufigkeiten der Einermengenprofile durch endliche
Zählung und erlauben, zwei verschiedene ausgezeichnete Elemente
vorzugeben. Der folgende Satz verbindet dies mit der besonderen Form des
Produktquadrats.

\Needspace{22\baselineskip}
\FormulaThmDeltaK[Rekonstruktion einer endlichen Halbgruppe mit zwei Produktwerten]{%
  \begin{aligned}[t]
    &\Finite{S}\dsep\SemigroupWithZeroAlg{S}{\star}{0_S}
      \dsep S^2=\{0_S,c\}\dsep c\neq0_S\\[-2pt]
    &\dsep\SemigroupAlg{T}{\diamond}
      \dsep\SemigroupIso{\Phi}
        {\PowNE{S},\star_{\mathcal P}}
        {\PowNE{T},\diamond_{\mathcal P}}\\[-2pt]
    &\vdash\exists f\,\SemigroupIso{f}{S,\star}{T,\diamond}
  \end{aligned}
}{FiniteTwoProductZeroReconstruction}{%
  \DeltaRow{Mengen}{S\dsep T\dsep0_S\dsep c}
  \DeltaRow{Funktionen}{\star\dsep\diamond\dsep\Phi\dsep f}
}

Dabei ist \(S^2=\{x\star y\mid x,y\in S\}\) das Produktquadrat,
und \(\PowNE{S}\) besteht ausschließlich aus den nichtleeren Teilmengen.
Endlichkeit und ein Nullelement von \(T\) brauchen nicht vorausgesetzt
zu werden: Beides wird im Beweis aus dem Potenzhalbgruppenisomorphismus
gewonnen.

\noindent\emph{Beweis:} Die \NullproductsReadingLink{} entwickelt den
vollständigen Beweis mit Beispielen. Die \NullproductsProofsLink{}
führen sämtliche dafür zusätzlich benötigten Ableitungen im
Lemmon-Stil aus. Der behauptete Isomorphismus wird abschließend hier
hergeleitet: \NullproductsProofLink{NPReconstructionConclusion}.

Der Satz behandelt diesen besonderen Fall der Rekonstruktionsfrage.
Er setzt weder eine Erhaltung der Einermengen voraus noch entscheidet er
die allgemeine Vermutung für beliebige endliche Halbgruppen.
